\chapter{Kế hoạch triển khai}
Phụ lục này trình bày kế hoạch triển khai đồ án theo 5 giai đoạn, trải dài từ tháng 2 đến tháng 9 năm 2026, với tháng 8 dành cho kiểm thử toàn diện và tháng 9 dành cho đóng gói và nghiệm thu.

\section{Biểu đồ Gantt}

\begin{figure}[H]
\centering
\begin{ganttchart}[
    hgrid,
    vgrid,
    x unit=1.05cm,
    y unit chart=0.7cm,
    bar height=0.55,
    bar/.append style={fill=blue!40},
    milestone/.append style={fill=red!70},
    title height=1,
    title/.style={fill=blue!10},
    group/.append style={fill=blue!60},
]{1}{8}
    \gantttitle{Kế hoạch triển khai DigitalBookLLM (02/2026 -- 09/2026)}{8} \\
    \gantttitlelist{Th.02,Th.03,Th.04,Th.05,Th.06,Th.07,Th.08,Th.09}{1} \\

    \ganttgroup{GĐ1: Khởi tạo \& yêu cầu}{1}{1} \\
    \ganttbar{Khảo sát vấn đề, đối thủ}{1}{1} \\
    \ganttmilestone{Chốt phạm vi đề tài}{1} \\

    \ganttgroup{GĐ2: Phân tích \& thiết kế}{2}{2} \\
    \ganttbar{Cơ sở lý thuyết, yêu cầu FR/NFR/UC}{2}{2} \\
    \ganttbar{Thiết kế kiến trúc, CSDL, UI/UX}{2}{2} \\

    \ganttgroup{GĐ3: Hiện thực hoá}{3}{6} \\
    \ganttbar{Backend: xác thực, workspace, hàng đợi}{3}{3} \\
    \ganttbar{Pipeline nạp tài liệu (OCR, chunking)}{3}{4} \\
    \ganttbar{Pipeline RAG \& bộ định tuyến LLM}{4}{5} \\
    \ganttbar{Frontend: trình đọc, chat, highlight}{4}{6} \\
    \ganttbar{Đồ thị tri thức cá nhân hoá}{5}{6} \\
    \ganttmilestone{Bản build đầu-cuối (MVP)}{6} \\

    \ganttgroup{GĐ4: Kiểm thử}{7}{7} \\
    \ganttbar{Kiểm thử đơn vị, tích hợp, bảo mật}{7}{7} \\
    \ganttbar{Kiểm thử hiệu năng \& đánh giá RAG}{7}{7} \\

    \ganttgroup{GĐ5: Đóng gói}{8}{8} \\
    \ganttbar{Triển khai công khai, hoàn thiện báo cáo}{8}{8} \\
    \ganttmilestone{Nghiệm thu đồ án}{8}
\end{ganttchart}
\caption{Biểu đồ Gantt kế hoạch triển khai đồ án}
\label{fig:gantt}
\end{figure}

\section{Diễn giải các giai đoạn}
\begin{longtable}{|C{0.10\textwidth}|C{0.16\textwidth}|L{0.62\textwidth}|}
\caption{Diễn giải chi tiết 5 giai đoạn triển khai}\\
\hline
\textbf{Giai đoạn} & \textbf{Thời gian} & \textbf{Công việc cụ thể} \\\hline
\endfirsthead
\hline
\textbf{Giai đoạn} & \textbf{Thời gian} & \textbf{Công việc cụ thể} \\\hline
\endhead
GĐ1 & Th.02/2026 & Khởi tạo dự án, khảo sát nỗi đau người dùng và đối thủ cạnh tranh, xác định mục tiêu và phạm vi đề tài. \\\hline
GĐ2 & Th.03/2026 & Xây dựng cơ sở lý thuyết (LLM, RAG, Vector DB, Knowledge Graph); phân tích yêu cầu (stakeholder, user journey, user story, FR/NFR, use case); thiết kế kiến trúc hệ thống, cơ sở dữ liệu và UI/UX. \\\hline
GĐ3 & Th.04--07/2026 & Hiện thực hoá theo từng phân hệ: xác thực \& workspace, pipeline nạp tài liệu, pipeline RAG \& bộ định tuyến LLM, giao diện đọc \& chat, đồ thị tri thức cá nhân hoá. Các phân hệ backend và frontend triển khai song song, tích hợp dần vào cuối giai đoạn. \\\hline
GĐ4 & Th.08/2026 & Kiểm thử toàn diện: đơn vị, tích hợp, bảo mật (IDOR), hiệu năng/tải đồng thời, đánh giá chất lượng truy xuất RAG (recall@k, TTFB). \\\hline
GĐ5 & Th.09/2026 & Triển khai công khai (Vercel/nền tảng container hoá + Supabase), hoàn thiện tài liệu báo cáo, chuẩn bị và thực hiện nghiệm thu đồ án. \\\hline
\end{longtable}

\textit{Ghi chú: mốc thời gian chi tiết theo tuần (phân bổ nhân lực, người phụ trách từng phân hệ) chưa được đưa vào bảng trên; xem tệp theo dõi công việc còn lại đính kèm đồ án để bổ sung theo tiến độ thực tế của nhóm.}
